\documentclass[12pt,a4paper]{report}
\usepackage[utf8]{inputenc}
\usepackage[T1]{fontenc}
\usepackage[french]{babel}
\usepackage{geometry}
\usepackage{graphicx}
\usepackage{hyperref}
\usepackage{booktabs}
\usepackage{longtable}
\usepackage{array}
\usepackage{tabularx}
\usepackage{xcolor}
\usepackage{colortbl}
\usepackage{enumitem}
\usepackage{fancyhdr}
\usepackage{titlesec}
\usepackage{tocloft}

\geometry{margin=2.5cm}

% Couleurs
\definecolor{emsiblue}{RGB}{0, 153, 51}
\definecolor{lightgray}{RGB}{240, 240, 240}
\definecolor{darkgray}{RGB}{100, 100, 100}

% En-têtes et pieds de page
\pagestyle{fancy}
\fancyhf{}
\fancyhead[L]{\small Cahier des Charges - Adhésion-App}
\fancyhead[R]{\small EMSI 2024-2025}
\fancyfoot[C]{\thepage}
\renewcommand{\headrulewidth}{0.4pt}

% Formatage des titres
\titleformat{\chapter}[display]
{\normalfont\huge\bfseries\color{emsiblue}}{\chaptertitlename\ \thechapter}{20pt}{\Huge}

\begin{document}

%==============================================================================
% PAGE DE GARDE
%==============================================================================
\begin{titlepage}
    \centering
    \vspace*{1cm}
    \begin{figure}
        \centering
        \includegraphics[width=0.5\linewidth]{EMSI.png}
        \label{fig:placeholder}
    \end{figure}
    
    {\Large \textbf{École Marocaine des Sciences de l'Ingénieur (EMSI)}}\\[0.5cm]
    {\large Filière : Ingénierie Informatique}\\[0.8cm]
    
    \rule{\linewidth}{0.5mm}\\[0.4cm]
    {\Huge \textbf{\color{emsiblue}CAHIER DES CHARGES}}\\[0.3cm]
    \rule{\linewidth}{0.5mm}\\[0.8cm]
    
    {\LARGE \textbf{Application de Prédiction des Comportements\\d'Adhésion aux Thérapies}}\\[0.5cm]
    {\Large \textbf{Basée sur Profils Psychométriques}}\\[0.3cm]
    {\large \textit{Adhésion-App / Serenity}}\\[1cm]
    
    \begin{tabular}{ll}
        \textbf{Projet :} & Projet Académique Transversal \\[0.2cm]
        \textbf{Matières :} & Management de la Qualité, Dev. Multiplateforme, JEE2, PFA \\[0.2cm]
        \textbf{Référence :} & CDC-AdhesionApp-v1.0 \\[0.2cm]
        \textbf{Version :} & 1.0 \\[0.2cm]
        \textbf{Date :} & 01 Novembre 2025 \\
    \end{tabular}\\[2cm]
    
    \textbf{Équipe de Projet :}\\[0.3cm]
    \begin{tabular}{|l|l|}
        \hline
        \rowcolor{emsiblue}\textcolor{white}{\textbf{Nom}} & \textcolor{white}{\textbf{Rôle}} \\
        \hline
        ISSAME Imad & Développeur / Testeur \\
        \hline
        AGOUMI Mohammed Amine & Développeur / Testeur \\
        \hline
        JABRANE Mohamed Yahya & Développeur / Testeur \\
        \hline
    \end{tabular}\\[2cm]
    
    \vfill
    {\small Année Universitaire 2025-2026}
\end{titlepage}

%==============================================================================
% TABLE DES MATIÈRES
%==============================================================================
\tableofcontents
\newpage

%==============================================================================
% HISTORIQUE DES VERSIONS
%==============================================================================
\chapter*{Historique des Versions}
\addcontentsline{toc}{chapter}{Historique des Versions}

\begin{tabularx}{\textwidth}{|c|c|X|l|}
    \hline
    \rowcolor{emsiblue}\textcolor{white}{\textbf{Version}} & \textcolor{white}{\textbf{Date}} & \textcolor{white}{\textbf{Description}} & \textcolor{white}{\textbf{Auteur}} \\
    \hline
    0.1 & 28/10/2025 & Version initiale - Définition du périmètre & Équipe \\
    \hline
    0.5 & 29/10/2025 & Ajout des spécifications fonctionnelles & Équipe \\
    \hline
    0.8 & 29/10/2025 & Spécifications techniques complètes & Équipe \\
    \hline
    1.0 & 01/11/2025 & Version finale validée & Équipe \\
    \hline
\end{tabularx}

\newpage

%==============================================================================
% CHAPITRE 1 : INTRODUCTION
%==============================================================================
\chapter{Introduction}

\section{Contexte du Projet}

Dans le cadre de notre formation à l'École Marocaine des Sciences de l'Ingénieur (EMSI), nous réalisons un projet académique transversal impliquant plusieurs matières :
\begin{itemize}
    \item \textbf{Management de la Qualité} : Mise en place de processus qualité et métriques
    \item \textbf{Développement Multiplateforme} : Application mobile Flutter
    \item \textbf{JEE2} : Backend Spring Boot avec architecture REST et Application web Angular
    \item \textbf{Projet de Fin d'Année (PFA)} : Intégration d'intelligence artificielle
\end{itemize}

L'adhésion thérapeutique représente un enjeu majeur de santé publique. Selon l'OMS, environ 50\% des patients atteints de maladies chroniques ne respectent pas leurs prescriptions médicales. Cette non-adhésion entraîne des complications de santé, une augmentation des coûts médicaux et une réduction de l'efficacité des traitements.

\section{Objectifs du Projet}

\subsection{Objectif Principal}
Développer une application multiplateforme permettant de :
\begin{itemize}
    \item \textbf{Évaluer} les profils psychométriques des patients via des tests validés cliniquement
    \item \textbf{Prédire} les comportements d'adhésion aux thérapies
    \item \textbf{Accompagner} les patients avec des recommandations personnalisées
    \item \textbf{Suivre} l'évolution de l'adhésion dans le temps
\end{itemize}

\subsection{Objectifs Spécifiques}
\begin{enumerate}
    \item Implémenter des tests psychométriques validés (MMAS-8, PHQ-9, GAD-7, BMQ)
    \item Développer un modèle de prédiction multi-factoriel basé sur la recherche
    \item Intégrer l'intelligence artificielle pour la personnalisation
    \item Fournir une interface utilisateur intuitive et accessible
    \item Assurer la qualité du code et la sécurité des données
\end{enumerate}

\section{Périmètre du Projet}

\subsection{Inclus dans le Périmètre}
\begin{itemize}
    \item Application mobile (Flutter) pour les patients et l'administrateur
    \item Application web Angular pour les patients
    \item Dashboard administrateur Angular
    \item API REST Spring Boot
    \item Module d'intelligence artificielle
    \item Base de données PostgreSQL
\end{itemize}

\subsection{Exclus du Périmètre}
\begin{itemize}
    \item Intégration avec des systèmes hospitaliers existants
    \item Application iOS native
    \item Téléconsultation en temps réel
\end{itemize}

\newpage

%==============================================================================
% CHAPITRE 2 : ANALYSE DES BESOINS
%==============================================================================
\chapter{Analyse des Besoins}

\section{Acteurs du Système}

\begin{tabularx}{\textwidth}{|l|X|}
    \hline
    \rowcolor{emsiblue}\textcolor{white}{\textbf{Acteur}} & \textcolor{white}{\textbf{Description}} \\
    \hline
    \textbf{Patient} & Utilisateur principal qui passe les tests, consulte ses prédictions et suit son traitement \\
    \hline
    \textbf{Administrateur} & Gère les utilisateurs, consulte les analytics et supervise le système \\
    \hline
    \textbf{Système IA} & Génère des recommandations, analyses et messages motivationnels \\
    \hline
\end{tabularx}

\section{Besoins Fonctionnels}

\subsection{Module Authentification (AUTH)}

\begin{longtable}{|p{2cm}|p{4cm}|p{7cm}|p{1.5cm}|}
    \hline
    \rowcolor{emsiblue}\textcolor{white}{\textbf{ID}} & \textcolor{white}{\textbf{Fonctionnalité}} & \textcolor{white}{\textbf{Description}} & \textcolor{white}{\textbf{Priorité}} \\
    \hline
    \endfirsthead
    \hline
    \rowcolor{emsiblue}\textcolor{white}{\textbf{ID}} & \textcolor{white}{\textbf{Fonctionnalité}} & \textcolor{white}{\textbf{Description}} & \textcolor{white}{\textbf{Priorité}} \\
    \hline
    \endhead
    
    AUTH-01 & Inscription & Le patient peut créer un compte avec email, mot de passe, nom, date de naissance et genre & Haute \\
    \hline
    AUTH-02 & Connexion & Authentification par email et mot de passe avec gestion de session & Haute \\
    \hline
    AUTH-03 & Déconnexion & Terminer la session utilisateur et nettoyer les données locales & Haute \\
    \hline
    AUTH-04 & Onboarding & Processus de tests initiaux obligatoires (minimum 2 tests) après inscription & Haute \\
    \hline
    AUTH-05 & Gestion Consentement & Recueil du consentement pour le partage de données & Moyenne \\
    \hline
\end{longtable}

\subsection{Module Tests Psychométriques (TEST)}

\begin{longtable}{|p{2cm}|p{4cm}|p{7cm}|p{1.5cm}|}
    \hline
    \rowcolor{emsiblue}\textcolor{white}{\textbf{ID}} & \textcolor{white}{\textbf{Fonctionnalité}} & \textcolor{white}{\textbf{Description}} & \textcolor{white}{\textbf{Priorité}} \\
    \hline
    \endfirsthead
    
    TEST-01 & Catalogue de Tests & Afficher la liste des tests disponibles avec descriptions & Haute \\
    \hline
    TEST-02 & MMAS-8 & Test de Morisky (8 items) pour mesurer l'adhésion médicamenteuse & Haute \\
    \hline
    TEST-03 & PHQ-9 & Questionnaire de santé du patient pour évaluer la dépression & Haute \\
    \hline
    TEST-04 & GAD-7 & Échelle d'anxiété généralisée (7 items) & Haute \\
    \hline
    TEST-05 & BMQ & Beliefs about Medicines Questionnaire - croyances sur les médicaments & Haute \\
    \hline
    TEST-06 & MARS-5 & Medication Adherence Report Scale - attitudes envers les médicaments & Moyenne \\
    \hline
    TEST-07 & Brief-IPQ & Perception de la maladie (version courte) & Moyenne \\
    \hline
    TEST-08 & SE-CHRONIC & Auto-efficacité pour maladies chroniques & Moyenne \\
    \hline
    TEST-09 & TSS & Therapeutic Social Support - support social thérapeutique & Basse \\
    \hline
    TEST-10 & Scoring Automatique & Calcul automatique des scores selon les grilles validées & Haute \\
    \hline
    TEST-11 & Interprétation & Génération d'interprétations cliniques basées sur les scores & Haute \\
    \hline
    TEST-12 & Historique & Consultation de l'historique des tests passés & Moyenne \\
    \hline
\end{longtable}

\subsection{Module Profil Psychologique (PROF)}

\begin{longtable}{|p{2cm}|p{4cm}|p{7cm}|p{1.5cm}|}
    \hline
    \rowcolor{emsiblue}\textcolor{white}{\textbf{ID}} & \textcolor{white}{\textbf{Fonctionnalité}} & \textcolor{white}{\textbf{Description}} & \textcolor{white}{\textbf{Priorité}} \\
    \hline
    
    PROF-01 & Génération Profil & Création automatique du profil psychologique après les tests & Haute \\
    \hline
    PROF-02 & Classification & Attribution d'un type de profil (Anxieux, Résilient, Motivé, etc.) & Haute \\
    \hline
    PROF-03 & Scores Dimensionnels & Calcul des scores : anxiété, dépression, motivation, auto-efficacité, support social & Haute \\
    \hline
    PROF-04 & Score de Risque & Calcul du score de risque de non-adhésion (0-1) & Haute \\
    \hline
    PROF-05 & Résumé IA & Génération d'un résumé personnalisé par l'IA & Moyenne \\
    \hline
    PROF-06 & Consultation Profil & Visualisation du profil avec graphiques & Moyenne \\
    \hline
\end{longtable}

\subsection{Module Prédiction d'Adhésion (PRED)}

\begin{longtable}{|p{2cm}|p{4cm}|p{7cm}|p{1.5cm}|}
    \hline
    \rowcolor{emsiblue}\textcolor{white}{\textbf{ID}} & \textcolor{white}{\textbf{Fonctionnalité}} & \textcolor{white}{\textbf{Description}} & \textcolor{white}{\textbf{Priorité}} \\
    \hline
    
    PRED-01 & Modèle Multi-factoriel & Prédiction basée sur 7 facteurs pondérés selon la recherche & Haute \\
    \hline
    PRED-02 & Facteurs de Risque & Identification et affichage des facteurs de risque personnels & Haute \\
    \hline
    PRED-03 & Niveau de Risque & Classification : Faible, Modéré, Élevé, Critique & Haute \\
    \hline
    PRED-04 & Score de Confiance & Indication de la fiabilité de la prédiction & Moyenne \\
    \hline
    PRED-05 & Recommandations IA & Génération de recommandations personnalisées & Haute \\
    \hline
    PRED-06 & Historique Prédictions & Suivi de l'évolution des prédictions dans le temps & Moyenne \\
    \hline
\end{longtable}

\textbf{Détail du Modèle de Prédiction :} \\

\begin{tabularx}{\textwidth}{|l|c|X|}
    \hline
    \rowcolor{emsiblue}\textcolor{white}{\textbf{Facteur}} & \textcolor{white}{\textbf{Poids}} & \textcolor{white}{\textbf{Source}} \\
    \hline
    Comportement Passé & 35\% & Historique des prises de doses (DoseLog) \\
    \hline
    Score MMAS-8 & 20\% & Auto-évaluation validée de l'adhésion \\
    \hline
    Dépression (PHQ-9) & 15\% & Impact majeur sur l'adhésion \\
    \hline
    Croyances (BMQ) & 10\% & Différentiel Nécessité-Préoccupations \\
    \hline
    Anxiété (GAD-7) & 8\% & Barrière potentielle \\
    \hline
    Complexité Traitement & 7\% & Nombre de médicaments et doses quotidiennes \\
    \hline
    Patterns Temporels & 5\% & Analyse des jours/heures problématiques \\
    \hline
\end{tabularx}

\subsection{Module Traitement (TREAT)}

\begin{longtable}{|p{2cm}|p{4cm}|p{7cm}|p{1.5cm}|}
    \hline
    \rowcolor{emsiblue}\textcolor{white}{\textbf{ID}} & \textcolor{white}{\textbf{Fonctionnalité}} & \textcolor{white}{\textbf{Description}} & \textcolor{white}{\textbf{Priorité}} \\
    \hline
    
    TREAT-01 & Création Plan & Génération d'un plan de traitement personnalisé & Haute \\
    \hline
    TREAT-02 & Génération IA & Contenu du plan généré par l'IA selon le profil & Haute \\
    \hline
    TREAT-03 & Médicaments & Gestion des médicaments (nom, dosage, fréquence, horaires) & Haute \\
    \hline
    TREAT-04 & Tâches Quotidiennes & Liste des tâches journalières à accomplir & Moyenne \\
    \hline
    TREAT-05 & Objectifs Hebdo & Définition et suivi des objectifs par semaine & Moyenne \\
    \hline
    TREAT-06 & Progression & Calcul et affichage du pourcentage de progression & Haute \\
    \hline
    TREAT-07 & Statuts Plan & Gestion des statuts : Brouillon, Actif, Pausé, Terminé, Annulé & Moyenne \\
    \hline
\end{longtable}

\subsection{Module Suivi des Doses (DOSE)}

\begin{longtable}{|p{2cm}|p{4cm}|p{7cm}|p{1.5cm}|}
    \hline
    \rowcolor{emsiblue}\textcolor{white}{\textbf{ID}} & \textcolor{white}{\textbf{Fonctionnalité}} & \textcolor{white}{\textbf{Description}} & \textcolor{white}{\textbf{Priorité}} \\
    \hline
    
    DOSE-01 & Doses du Jour & Affichage des doses planifiées pour aujourd'hui & Haute \\
    \hline
    DOSE-02 & Marquage Prise & Enregistrer qu'une dose a été prise avec horodatage & Haute \\
    \hline
    DOSE-03 & Marquage Oubli & Enregistrer qu'une dose a été oubliée avec raison optionnelle & Haute \\
    \hline
    DOSE-04 & Raisons d'Oubli & Liste prédéfinie : Oublié, Effets secondaires, Pas disponible, etc. & Moyenne \\
    \hline
    DOSE-05 & Délai de Prise & Calcul du retard/avance par rapport à l'horaire prévu & Moyenne \\
    \hline
    DOSE-06 & Notes & Ajout de notes personnelles (ressenti, effets) & Basse \\
    \hline
    DOSE-07 & Streak & Calcul des séries de jours consécutifs d'adhésion & Moyenne \\
    \hline
    DOSE-08 & Statistiques & Taux d'adhésion global et par médicament & Haute \\
    \hline
\end{longtable}

\subsection{Module Recommandations (REC)}

\begin{longtable}{|p{2cm}|p{4cm}|p{7cm}|p{1.5cm}|}
    \hline
    \rowcolor{emsiblue}\textcolor{white}{\textbf{ID}} & \textcolor{white}{\textbf{Fonctionnalité}} & \textcolor{white}{\textbf{Description}} & \textcolor{white}{\textbf{Priorité}} \\
    \hline
    
    REC-01 & Génération & Création de recommandations basées sur le profil & Haute \\
    \hline
    REC-02 & Catégorisation & Catégories : Médicaments, Style de vie, Santé mentale, Social, Exercice & Haute \\
    \hline
    REC-03 & Priorité & Attribution de priorité (1-5) à chaque recommandation & Moyenne \\
    \hline
    REC-04 & Étapes Actions & Liste d'étapes concrètes pour chaque recommandation & Haute \\
    \hline
    REC-05 & Bénéfices & Explication des bénéfices attendus & Moyenne \\
    \hline
    REC-06 & Difficulté & Niveau de difficulté : Facile, Moyen, Difficile & Basse \\
    \hline
    REC-07 & Suivi Complétion & Marquage des recommandations complétées & Moyenne \\
    \hline
\end{longtable}

\subsection{Module Motivation IA (MOTIV)}

\begin{longtable}{|p{2cm}|p{4cm}|p{7cm}|p{1.5cm}|}
    \hline
    \rowcolor{emsiblue}\textcolor{white}{\textbf{ID}} & \textcolor{white}{\textbf{Fonctionnalité}} & \textcolor{white}{\textbf{Description}} & \textcolor{white}{\textbf{Priorité}} \\
    \hline
    
    MOTIV-01 & Message Motivation & Génération de messages motivationnels personnalisés & Haute \\
    \hline
    MOTIV-02 & Chat IA & Conversation avec l'assistant IA pour support psychologique & Haute \\
    \hline
    MOTIV-03 & Historique Chat & Conservation de l'historique des conversations & Moyenne \\
    \hline
    MOTIV-04 & Contexte Patient & Utilisation du profil et score d'adhésion dans les réponses & Haute \\
    \hline
\end{longtable}

\subsection{Module Émotions (EMOT)}

\begin{longtable}{|p{2cm}|p{4cm}|p{7cm}|p{1.5cm}|}
    \hline
    \rowcolor{emsiblue}\textcolor{white}{\textbf{ID}} & \textcolor{white}{\textbf{Fonctionnalité}} & \textcolor{white}{\textbf{Description}} & \textcolor{white}{\textbf{Priorité}} \\
    \hline
    
    EMOT-01 & Détection Faciale & Reconnaissance des émotions via caméra (8 émotions) & Moyenne \\
    \hline
    EMOT-02 & Analyse Texte & Détection d'émotions dans le texte écrit & Basse \\
    \hline
    EMOT-03 & Recommandations & Suggestions basées sur l'émotion détectée & Moyenne \\
    \hline
\end{longtable}

\subsection{Module Administration (ADMIN)}

\begin{longtable}{|p{2cm}|p{4cm}|p{7cm}|p{1.5cm}|}
    \hline
    \rowcolor{emsiblue}\textcolor{white}{\textbf{ID}} & \textcolor{white}{\textbf{Fonctionnalité}} & \textcolor{white}{\textbf{Description}} & \textcolor{white}{\textbf{Priorité}} \\
    \hline
    
    ADMIN-01 & Dashboard & Tableau de bord avec statistiques globales & Haute \\
    \hline
    ADMIN-02 & Liste Patients & Affichage et recherche des patients inscrits & Haute \\
    \hline
    ADMIN-03 & Détail Patient & Consultation du profil complet d'un patient & Haute \\
    \hline
    ADMIN-04 & Analytics & Graphiques et métriques d'utilisation & Moyenne \\
    \hline
    ADMIN-05 & Gestion Tests & Administration des définitions de tests & Moyenne \\
    \hline
    ADMIN-06 & Export Données & Export des données anonymisées pour analyse & Basse \\
    \hline
\end{longtable}

\section{Besoins Non Fonctionnels}

\begin{longtable}{|p{2.5cm}|p{4cm}|p{7cm}|}
    \hline
    \rowcolor{emsiblue}\textcolor{white}{\textbf{Catégorie}} & \textcolor{white}{\textbf{Exigence}} & \textcolor{white}{\textbf{Description}} \\
    \hline
    \endfirsthead
    
    \multirow{\textbf{Performance}} 
    & Temps de réponse API & < 500ms pour 95\% des requêtes \\
    \cline{2-3}
    & Charge simultanée & Support de 100 utilisateurs concurrents \\
    \cline{2-3}
    & Disponibilité & 99\% de disponibilité hors maintenance \\
    \hline
    
    \multirow{\textbf{Sécurité}}
    & Chiffrement MDP & BCrypt avec salt \\
    \cline{2-3}
    & Données sensibles & Stockage sécurisé des données médicales \\
    \cline{2-3}
    & Sessions & Gestion sécurisée avec expiration \\
    \hline
    
    \multirow{\textbf{Portabilité}}
    & Mobile & Android (API 21+), iOS (via Flutter) \\
    \cline{2-3}
    & Web & Navigateurs modernes (Chrome, Firefox, Edge) \\
    \hline
    
    \multirow{\textbf{Maintenabilité}}
    & Architecture & Modulaire par domaine (DDD) \\
    \cline{2-3}
    & Documentation & Code documenté, API Swagger \\
    \hline
    
    \multirow{\textbf{Utilisabilité}}
    & Interface & Design moderne et intuitif \\
    \cline{2-3}
    & Accessibilité & Support des lecteurs d'écran \\
    \hline
\end{longtable}

\newpage

%==============================================================================
% CHAPITRE 3 : ARCHITECTURE TECHNIQUE
%==============================================================================
\chapter{Architecture Technique}

\section{Vue d'Ensemble}

\begin{center}
\begin{tabularx}{\textwidth}{|l|X|}
    \hline
    \rowcolor{emsiblue}\textcolor{white}{\textbf{Composant}} & \textcolor{white}{\textbf{Technologie}} \\
    \hline
    Backend API & Spring Boot 3.5.7 (Java 17) \\
    \hline
    Base de données & PostgreSQL 15 \\
    \hline
    Application Mobile & Flutter (Dart 3.9.2) \\
    \hline
    Application Web Patient & Angular 20 \\
    \hline
    Dashboard Admin & Angular 20 \\
    \hline
    Backend IA & Flask (Python) \\
    \hline
    API IA & OpenRouter (Tongyi/Gemini) \\
    \hline
    Détection Émotions & EmotiEffLib (ONNX Runtime) \\
    \hline
    Conteneurisation & Docker, Docker Compose \\
    \hline
    Qualité Code & SonarQube \\
    \hline
\end{tabularx}
\end{center}

\section{Architecture Backend}

L'architecture suit un pattern modulaire par domaine (Domain-Driven Design) :

\begin{verbatim}
src/main/java/com/projet/adhesionapp/
├── identity/       # Gestion utilisateurs, authentification
├── profile/        # Profils psychologiques
├── assessment/     # Tests psychométriques
├── treatment/      # Plans de traitement
├── analytics/      # Prédictions d'adhésion
├── recommendation/ # Recommandations personnalisées
├── habit/          # Suivi des doses
├── ai/             # Intégration IA
├── common/         # Utilitaires partagés
└── config/         # Configuration Spring
\end{verbatim}

\section{Modèle de Données}

\subsection{Entités Principales}

\begin{center}
\begin{tabularx}{\textwidth}{|l|X|}
    \hline
    \rowcolor{emsiblue}\textcolor{white}{\textbf{Entité}} & \textcolor{white}{\textbf{Description}} \\
    \hline
    User & Utilisateur (patient) avec informations personnelles \\
    \hline
    PsychologicalProfile & Profil psychologique calculé \\
    \hline
    TestDefinition & Définition d'un test psychométrique \\
    \hline
    TestSession & Session de passage d'un test \\
    \hline
    Answer & Réponse à une question de test \\
    \hline
    TreatmentPlan & Plan de traitement personnalisé \\
    \hline
    Medication & Médicament prescrit \\
    \hline
    DailyTask & Tâche quotidienne du plan \\
    \hline
    DoseLog & Log de prise/oubli de dose \\
    \hline
    Prediction & Prédiction d'adhésion \\
    \hline
    Recommendation & Recommandation personnalisée \\
    \hline
\end{tabularx}
\end{center}

\section{API REST}

\subsection{Description des Services API}

L'API REST expose plusieurs catégories de services permettant la communication entre les applications clientes et le backend : \\

\textbf{Services d'Authentification :} Ces services permettent la gestion du cycle de vie des utilisateurs. Un service POST gère l'inscription des nouveaux patients en créant leur compte avec les informations personnelles. Un autre service POST assure la connexion en vérifiant les identifiants et en établissant une session sécurisée. \\

\textbf{Services de Tests Psychométriques :} Un service GET permet de récupérer la liste complète des tests disponibles (MMAS-8, PHQ-9, GAD-7, etc.) avec leurs descriptions. Un service POST permet de démarrer une nouvelle session de test pour un patient. Un autre service POST permet de soumettre les réponses d'une session de test et déclenche le calcul automatique des scores. \\

\textbf{Services de Profils Psychologiques :} Un service GET permet de récupérer le profil psychologique complet d'un utilisateur incluant ses scores dimensionnels et son type de profil. Un service POST permet de créer ou régénérer un profil basé sur les derniers résultats de tests. \\

\textbf{Services de Prédiction :} Un service GET génère une prédiction d'adhésion pour un utilisateur en appliquant le modèle multi-factoriel. Un autre service GET calcule et retourne le score d'adhésion actuel basé sur l'historique des prises de doses. \\

\textbf{Services de Traitements :} Un service POST permet de créer un nouveau plan de traitement personnalisé avec génération de contenu par l'IA. Un service GET récupère tous les plans de traitement associés à un utilisateur. \\

\textbf{Services de Suivi des Doses :} Un service GET retourne la liste des doses planifiées pour la journée en cours. Deux services POST permettent respectivement de marquer une dose comme prise (avec horodatage) ou comme oubliée (avec possibilité d'indiquer une raison). \\

\textbf{Services d'Intelligence Artificielle :} Un service POST génère des messages motivationnels personnalisés selon le profil et le contexte du patient. Un autre service POST permet d'avoir une conversation interactive avec l'assistant IA pour un support psychologique continu. \\

\newpage

%==============================================================================
% CHAPITRE 4 : INTERFACES UTILISATEUR
%==============================================================================
\chapter{Interfaces Utilisateur}

\section{Application Mobile (Flutter)}

\subsection{Écrans Principaux}

\begin{longtable}{|l|p{8cm}|}
    \hline
    \rowcolor{emsiblue}\textcolor{white}{\textbf{Écran}} & \textcolor{white}{\textbf{Description}} \\
    \hline
    
    Login & Formulaire de connexion avec email/mot de passe \\
    \hline
    Register & Inscription avec validation des champs \\
    \hline
    Onboarding & Passage des tests initiaux obligatoires \\
    \hline
    Dashboard & Vue d'ensemble avec score wellness, statistiques, raccourcis \\
    \hline
    Tests & Catalogue et passage des tests psychométriques \\
    \hline
    Profil & Visualisation du profil psychologique avec graphiques \\
    \hline
    Prédictions & Affichage de la prédiction d'adhésion et facteurs de risque \\
    \hline
    Doses du jour & Liste des doses avec boutons de confirmation \\
    \hline
    Statistiques & Graphiques d'adhésion sur différentes périodes \\
    \hline
    Médicaments & Gestion des médicaments personnels \\
    \hline
    Plan de Traitement & Détail du plan avec progression \\
    \hline
    Recommandations & Liste des recommandations personnalisées \\
    \hline
    Motivation & Chat avec l'assistant IA \\
    \hline
    Émotions & Détection faciale d'émotions via caméra \\
    \hline
    Historique & Historique des tests et prédictions \\
    \hline
\end{longtable}

\section{Application Web Patient (Angular)}

\begin{longtable}{|l|p{8cm}|}
    \hline
    \rowcolor{emsiblue}\textcolor{white}{\textbf{Vue}} & \textcolor{white}{\textbf{Fonctionnalités}} \\
    \hline
    
    Login / Register & Authentification et inscription des patients \\
    \hline
    Dashboard & Tableau de bord personnel avec score wellness et statistiques \\
    \hline
    Tests & Catalogue et passage des tests psychométriques en ligne \\
    \hline
    Profil & Visualisation et édition du profil psychologique \\
    \hline
    Prédictions & Consultation des prédictions d'adhésion et facteurs de risque \\
    \hline
    Médicaments & Gestion des médicaments et suivi des prises \\
    \hline
    Plans de Traitement & Consultation des plans de traitement actifs \\
    \hline
    Recommandations & Affichage des recommandations personnalisées \\
    \hline
    Motivation & Interface de chat avec l'assistant IA \\
    \hline
    Émotions & Détection d'émotions via webcam \\
    \hline
    Historique & Historique complet des tests et résultats \\
    \hline
\end{longtable}

\section{Dashboard Administrateur (Angular)}

\begin{longtable}{|l|p{8cm}|}
    \hline
    \rowcolor{emsiblue}\textcolor{white}{\textbf{Vue}} & \textcolor{white}{\textbf{Fonctionnalités}} \\
    \hline
    
    Login Admin & Authentification sécurisée pour administrateurs \\
    \hline
    Dashboard & KPIs globaux, graphiques d'utilisation, alertes \\
    \hline
    Gestion Patients & Liste paginée, recherche, filtres, détail patient \\
    \hline
    Données Utilisateurs & Consultation des profils et données psychologiques \\
    \hline
    Gestion Tests & Administration des définitions de tests \\
    \hline
    Analytics & Métriques avancées et exports de données \\
    \hline
\end{longtable}

\newpage

%==============================================================================
% CHAPITRE 5 : CONTRAINTES ET LIVRABLES
%==============================================================================
\chapter{Contraintes et Livrables}

\section{Contraintes du Projet}

\subsection{Contraintes Temporelles}
\begin{itemize}
    \item \textbf{Durée :} 8 semaines (01/11/2025 - 22/12/2025)
    \item \textbf{Deadline :} 22 Décembre 2025
\end{itemize}

\subsection{Contraintes Techniques}
\begin{itemize}
    \item Utilisation obligatoire de Spring Boot pour le backend
    \item Développement multiplateforme avec Flutter
    \item Base de données relationnelle PostgreSQL
    \item Respect des bonnes pratiques de développement
\end{itemize}

\subsection{Contraintes Qualité}

\subsubsection{Tests Unitaires (JUnit)}
\begin{itemize}
    \item Couverture de code minimale : \textbf{80\%}
    \item Intégration de JUnit dans le processus CI/CD pour exécution automatique
    \item Tests des services, contrôleurs et utilitaires
    \item Utilisation de Mockito pour les tests isolés
\end{itemize}

\subsubsection{Tests Fonctionnels (Selenium)}
\begin{itemize}
    \item Automatisation des tests fonctionnels de l'application web
    \item Focus sur les scénarios critiques (authentification, tests psychométriques, suivi des doses)
    \item Exécution des tests sur plusieurs navigateurs (Chrome, Firefox, Edge)
    \item Tests de bout en bout des parcours utilisateur
\end{itemize}

\subsubsection{Analyse de Qualité de Code (SonarQube)}
\begin{itemize}
    \item Analyse statique du code pour détecter les problèmes de qualité
    \item Seuils à ne pas dépasser :
    \begin{itemize}
        \item Bugs : 0 critique, < 10 majeurs
        \item Vulnérabilités : 0 critique, 0 haute
        \item Code Smells : < 100
        \item Duplication de code : < 5\%
    \end{itemize}
    \item Intégration CI/CD : SonarScanner exécuté automatiquement après chaque build
    \item Correction obligatoire des problèmes critiques avant merge
\end{itemize}

\subsubsection{Tests de Performance (JMeter)}
\begin{itemize}
    \item Développement de scénarios de test pour les cas critiques
    \item Configuration d'outils de monitoring pour surveiller les performances serveur
    \item Documentation complète incluant :
    \begin{itemize}
        \item Description des scénarios de test
        \item Résultats et métriques collectées
        \item Recommandations d'optimisation
    \end{itemize}
    \item Seuils de performance : temps de réponse < 500ms pour 95\% des requêtes
\end{itemize}

\subsubsection{Normes de Codage}
\begin{itemize}
    \item Respect des conventions Java (Google Java Style Guide)
    \item Respect des conventions Dart (Effective Dart)
    \item Documentation Javadoc/DartDoc pour les API publiques
\end{itemize}

\section{Livrables}

\begin{tabularx}{\textwidth}{|c|l|X|}
    \hline
    \rowcolor{emsiblue}\textcolor{white}{\textbf{N°}} & \textcolor{white}{\textbf{Livrable}} & \textcolor{white}{\textbf{Description}} \\
    \hline
    L1 & Cahier des Charges & Ce document \\
    \hline
    L2 & Plan Assurance Qualité & Processus et critères qualité\\
    \hline
    L3 & Code Source Backend & API Spring Boot complète \\
    \hline
    L4 & Code Source Mobile & Application Flutter\\
    \hline
    L5 & Code Source Web & Applications Angular \\
    \hline
    L6 & Rapport SonarQube & Analyse qualité du code \\
    \hline
    L7 & Rapport Tests & Résultats JUnit, JMeter, Selenium \\
    \hline
    L8 & Documentation API & Swagger/OpenAPI \\
    \hline
    L9 & Manuel Utilisateur & Guide d'utilisation \\
    \hline
\end{tabularx}

\section{Critères d'Acceptation}

\begin{enumerate}
    \item Toutes les fonctionnalités de priorité "Haute" sont implémentées
    \item Les tests unitaires passent avec succès
    \item Les tests de performance respectent les seuils définis
    \item Le code respecte les métriques SonarQube
    \item L'application est déployable via Docker
    \item La documentation est complète et à jour
\end{enumerate}

\newpage

%==============================================================================
% CHAPITRE 6 : GLOSSAIRE
%==============================================================================
\chapter{Glossaire}

\begin{tabularx}{\textwidth}{|l|X|}
    \hline
    \rowcolor{emsiblue}\textcolor{white}{\textbf{Terme}} & \textcolor{white}{\textbf{Définition}} \\
    \hline
    Adhésion thérapeutique & Degré auquel le patient suit les prescriptions médicales \\
    \hline
    MMAS-8 & Morisky Medication Adherence Scale - échelle validée d'adhésion \\
    \hline
    PHQ-9 & Patient Health Questionnaire - évaluation de la dépression \\
    \hline
    GAD-7 & Generalized Anxiety Disorder - échelle d'anxiété \\
    \hline
    BMQ & Beliefs about Medicines Questionnaire - croyances sur les médicaments \\
    \hline
    Profil psychométrique & Ensemble des scores aux tests psychologiques d'un patient \\
    \hline
    Score de risque & Probabilité calculée de non-adhésion (0-1) \\
    \hline
    DoseLog & Enregistrement d'un événement de prise de médicament \\
    \hline
    Streak & Série de jours consécutifs d'adhésion parfaite \\
    \hline
    OpenRouter & Service d'API pour accéder à différents modèles d'IA \\
    \hline
    EmotiEffLib & Bibliothèque de reconnaissance d'émotions faciales \\
    \hline
    ONNX & Open Neural Network Exchange - format de modèle ML portable \\
    \hline
\end{tabularx}

\newpage

%==============================================================================
% SIGNATURES
%==============================================================================
\chapter*{Validation du Document}
\addcontentsline{toc}{chapter}{Validation du Document}

\vspace{1cm}

\begin{tabularx}{\textwidth}{|X|X|}
    \hline
    \rowcolor{emsiblue}\textcolor{white}{\textbf{Rôle}} & \textcolor{white}{\textbf{Nom Complet}} \\
    \hline
    & & \\
    Chef de Projet & ISSAME Imad \\
    & & \\
    \hline
    & & \\
    Développeur/Testeur & AGOUMI Mohammed Amine \\
    & & \\
    \hline
    & & \\
    Développeur/Testeur & JABRANE Mohamed Yahya \\
    & & \\
    \hline
    & & \\
    Encadrant & M. ESSABAR Driss \\
    & & \\
    \hline
\end{tabularx}

\vspace{2cm}

\begin{center}
\textit{Document généré le 01 Novembre 2025}
\end{center}

\end{document}
